% TOP_PROBLEM: {"id":30007569,"problem_number":"LOCAL-30007569","title":"Manipulation-resistant learning","category_id":21,"category":{"id":21,"name":"economics","display_name":"Economics","description":"Open economic theory statements, published research agendas, and proposed scoped specifications; question provenance and historical priority are qualified per record.","slug":"economics","order_index":21,"created_at":"2026-10-08T00:00:00Z"},"status":"research_question","external_url":null,"legacy_ids":[],"published":false,"statement":"Characterize learning-versus-extraction guarantees for two privately informed persistent buyers facing an adaptive seller, extending the source’s solved single-agent setting.","economics":{"original_id":58,"field":"Game theory and strategic interaction","kind":"design","formalization_basis":"adapted_research_question","source_status":"research_agenda","priority_status":"earliest_found","priority_note":"This is the earliest source in which the relevant agenda was verified during this audit; absolute priority has not been established.","earliest_verified_reference":{"authors":["Kyohei Okumura"],"title":"Robust Learning with Private Information","year":2026,"url":"https://arxiv.org/pdf/2505.05341","doi":null,"locator":"§8 Conclusion, printed p.20, final paragraph (version3, 10February2026)","evidence_summary":"After solving its persistent single-agent screening model, the paper explicitly calls for richer markets with multiple learners and other payoff and information structures."},"current_reference":{"authors":["Kyohei Okumura"],"title":"Robust Learning with Private Information","year":2026,"url":"https://arxiv.org/pdf/2505.05341","doi":null,"locator":"§8 Conclusion, printed p.20, final paragraph (version3, 10February2026)","evidence_summary":"After solving its persistent single-agent screening model, the paper explicitly calls for richer markets with multiple learners and other payoff and information structures."},"formalization_note":"The exact two-buyer frontier is proposed; the verified conclusion, not the paper’s solved single-agent question, is the open-agenda source."},"statusQualification":"Published question or research agenda verified at the cited locator. The mathematical specification is adapted for this catalogue and includes additional assumptions; source publication does not establish first-ever priority or certify this exact model remains unsolved. The exact two-buyer frontier is proposed; the verified conclusion, not the paper’s solved single-agent question, is the open-agenda source.","statusReviewedAt":"2026-10-08"}
% FIRST_POSTED: 2026-10-08
% ENTRY_KIND: statement_only
% !TeX program = lualatex
\documentclass[11pt]{article}
\usepackage{fontspec}
\usepackage[a4paper,margin=25mm]{geometry}
\usepackage{amsmath,amssymb,mathtools}
\usepackage{xurl}
\usepackage[hidelinks]{hyperref}
\newcommand{\sref}[2]{\hyperlink{source:#1:#2}{[S#2]}}
\setlength{\emergencystretch}{3em}
\begin{document}
% problemId:30007569
\section*{Manipulation-resistant learning}\label{30007569}
\subsection{Definitions and mathematical statement}
Consider a repeated market with one seller and two buyers. Buyer \(i\) privately observes a persistent type \(\theta_i\) in a known finite set \(\Theta_i\subseteq[0,1]\); the seller knows a full-support joint prior \(F\) from a specified class \(\mathcal P\), but not realizations. Each period buyers choose finite actions, and a seller policy maps their actions into a feasible allocation of one good and nonnegative payments. Buyers observe only their own realized allocations and payments and may opt out for zero payoff. Before interaction each buyer commits to a learning algorithm \(L_i\). Define stationary learning loss \(R_T(L)\) against the best fixed action when seller policies and other buyers’ actions are independently drawn each period from fixed unknown distributions. Against adaptive sellers, let \(\Sigma^*(L,F)\) be revenue-maximizing seller responses to the committed algorithm profile. For a static revenue-optimal Bayesian mechanism \(M_F^*\), define extraction loss
\[
 S_T(L)=\sup_{F\in\mathcal P}\sup_{\sigma\in\Sigma^*(L,F)}\max_{i,\theta_i}
 \left(U_i(M_F^*\mid\theta_i)-E_{L,\sigma}[T^{-1}\textstyle\sum_t u_{it}\mid\theta_i]\right).
\]
Characterize the attainable tradeoff between sublinear \(R_T\) and vanishing or bounded \(S_T\), including necessary information and participation assumptions. Define buyer utility \(u_{it}=\theta_i x_{it}-p_{it}\) and ensure the static mechanism uses identical feasibility constraints. The cited paper solves a single-agent model and explicitly identifies multiple learners as a future direction; this two-buyer frontier is an adaptation and may require a different benchmark.

\par\textbf{Formulation and scope.} The exact two-buyer frontier is proposed; the verified conclusion, not the paper’s solved single-agent question, is the open-agenda source.

\par\textbf{Open-question provenance.} An explicit question or research agenda was verified at the source locator. This does not by itself certify that every later version remains unresolved \sref{30007569}{1}.

\par\textbf{Historical priority.} This is the earliest source in which the relevant agenda was verified during this audit; absolute priority has not been established.

\subsection{Short English statement}
Characterize learning-versus-extraction guarantees for two privately informed persistent buyers facing an adaptive seller, extending the source’s solved single-agent setting.

\subsection{Sources}
\begin{itemize}\raggedright
\item[\textnormal{[S1]}]\hypertarget{source:30007569:1}{} Kyohei Okumura. 2026. Robust Learning with Private Information. §8 Conclusion, printed p.20, final paragraph (version3, 10February2026).
\url{https://arxiv.org/pdf/2505.05341}.
{\small\textit{Earliest verified question/agenda reference; historical priority qualified below. After solving its persistent single-agent screening model, the paper explicitly calls for richer markets with multiple learners and other payoff and information structures.}}
\end{itemize}
\end{document}
